\documentclass{article}
\usepackage{amsmath,amssymb,hyperref}
\title{Pullback comparison direction: a geometric research checkpoint}
\begin{document}\maketitle
For a bifunctor $p$ and two cospans $X_i\to Z_i\leftarrow Y_i$, functoriality and the pullback universal property give
\[
p(X_1\times_{Z_1}Y_1,X_2\times_{Z_2}Y_2)\longrightarrow
p(X_1,X_2)\times_{p(Z_1,Z_2)}p(Y_1,Y_2).
\]
This direction is local parallel pullbacks to the global pullback. A colax comparison for a display-based crossing, with forward native context maps, may require the opposite direction. It cannot be obtained by relabelling lax and colax.

A reverse cone factorization is not automatic. In abelian groups let $F=(-)\otimes_{\mathbb Z}\mathbb Z/2$. The pullback of $\mathbb Z\xrightarrow{2}\mathbb Z\leftarrow0$ is $0$. After $F$, its pullback is $\mathbb Z/2$, and the canonical comparison is $0\to\mathbb Z/2$. No reverse map preserves the first cone projection, since that would factor the identity of $\mathbb Z/2$ through $0$. This tests one variable of a genuine tensor bifunctor, but is not asserted to instantiate all inputs of the relative crossing theorem.

Grandis--Par\'e, \emph{Intercategories}, Section 6.1, \url{https://www.mscs.dal.ca/~pare/IntercatIIsf.pdf}, derives noninvertible comparison from compatible chosen pullbacks; Section 6.4 gives spans of cospans. Identifying its axes and cell orientation with the present relative construction remains a theorem to prove. Possible routes are a differently oriented interaction, an actual adjunction/mate transfer, or independent exactness assumptions. No general impossibility of concurrency or feedback follows from this direction check. This direction check has received independent written-proof review.
\end{document}
